\section[Cote 44 : l'intégralité des nombres de Cartan à partir des axiomes]{Cote 44 : l'intégralité des nombres de Cartan à partir des axiomes\protect\verifie{Folder44}}\label{sec:44}

La cote 44, \emph{Jacobiennes intermédiaires : notes manuscrites (s.d.), lettre (1964)}, compte huit pages. Trois feuillets d'un programme sur les jacobiennes intermédiaires y entourent la lettre à Serre du 12 août 1964. La page~4 est étrangère au programme : c'est un feuillet barré d'une longue diagonale, qui énumère deux structures axiomatiques, un système de racines réduit et ce qu'on appelle aujourd'hui une donnée radicielle. La première est un couple $(P, \Delta)$, $P$ abélien libre de rang fini et $\Delta \subset P$ finie, soumis à trois conditions : a) $\Delta$ engendre $P$ ; b) si $a \in \Delta$, alors $-a \in \Delta$, \q{mais $na \notin \Delta$ si $n \neq 1, -1$, $n \in \mathbb Q$} ; c) pour tout $a \in \Delta$, il existe une \q{symétrie} $S_a$ qui envoie $a$ sur $-a$ et vérifie $S_a(\Delta) \subset \Delta$. Une note réalise toutes les symétries par la forme invariante $\sum_{a \in \Delta} \langle x, a\rangle\langle y, a\rangle$, écrit $S_\alpha x = x - \langle\gamma_\alpha, x\rangle\,\alpha$ avec $\gamma_\alpha \in P'^{\mathbb Q}$, puis, après un mot illisible, \q{que l'on a en fait $\gamma_\alpha \in P'$}, c'est-à-dire que les nombres de Cartan $c_{\alpha\beta} = \langle\gamma_\alpha, \beta\rangle$ sont entiers. Un crochet en donne la raison : chaque $\alpha$ serait indivisible dans $P$, parce qu'il peut faire partie d'un \q{système fondamental de racines} (\lot{44}{1}{4}). Le mot qui gouverne l'assertion et une bonne part du crochet sont illisibles. La lecture modernisée y voit l'intégralité tirée des axiomes comme une conséquence ; la piste de l'onglet Findings fait la réserve que les trois axiomes seuls ne la donnent pas.

\noindent\emph{Conventions.} Soient $V$ un espace vectoriel sur $\mathbb Q$ et $P$ un sous-groupe de $V$. Un élément $\alpha \in P$ est \emph{indivisible dans $P$} si, chaque fois que $\alpha = m y$ avec $m \in \mathbb Z$ et $y \in P$, on a $m = \pm 1$. Dans l'exemple, $P = \mathbb Z^2$, de base $(e_1, e_2)$, et les symétries sont calculées dans $\mathbb Q^2$.

\begin{enonce}[Page 4, de l'indivisibilité à l'intégralité\piste{44-cartan-integrality-from-axioms}]
Soient $\alpha \in P$ indivisible dans $P$ et $f : V \to \mathbb Q$ une application telle que $x - f(x)\,\alpha \in P$ pour tout $x \in P$. Alors $f(x) \in \mathbb Z$ pour tout $x \in P$. En particulier, si la symétrie $S_\alpha x = x - \langle\gamma_\alpha, x\rangle\,\alpha$ envoie $P$ dans $P$ et si $\alpha$ est indivisible dans $P$, les nombres $\langle\gamma_\alpha, \beta\rangle$, $\beta \in P$, sont entiers.
\end{enonce}

\begin{proof}
Soient $x \in P$ et $q = f(x)$ ; on a $q\alpha = x - (x - q\alpha) \in P$. Écrivons $q = a/b$ avec $b \geq 1$ et $a$, $b$ premiers entre eux, et choisissons $u, v \in \mathbb Z$ tels que $ua + vb = 1$. L'élément $y = u\,(q\alpha) + v\,\alpha$ est dans $P$, et $b y = u a\,\alpha + v b\,\alpha = \alpha$. Par indivisibilité, $b = \pm 1$, donc $b = 1$ et $q = a \in \mathbb Z$.
\end{proof}

\begin{remarque}
L'argument n'utilise ni la liberté de $P$, ni la finitude de $\Delta$, ni la forme invariante : il suffit que $S_\alpha$ préserve $P$ et soit de la forme $x \mapsto x - f(x)\,\alpha$. C'est le pas que le crochet attribue à l'indivisibilité.
\end{remarque}

\begin{contrexemple}[Page 4, les trois axiomes seuls\piste{44-cartan-integrality-from-axioms}]
Dans $P = \mathbb Z^2$, soit
\[ \Delta = \{\pm 3e_1, \pm 3e_2, \pm(e_1 + e_2), \pm(e_1 - e_2)\}. \]
\begin{enumerate}
\item $\Delta$ engendre $\mathbb Z^2$ (axiome a).
\item $\Delta = -\Delta$, et aucune racine n'est un multiple rationnel $q a$ d'une racine $a$ avec $q \neq \pm 1$ (axiome b, avec $n \in \mathbb Q$ comme sur la page).
\item Pour chaque $a \in \Delta$, il existe un automorphisme $S$ de $\mathbb Z^2$ tel que $S a = -a$ et $S(\Delta) \subset \Delta$ (axiome c).
\item Soit $f : \mathbb Z^2 \to \mathbb Q$ additive, avec $f(3e_1) = 2$, telle que $x \mapsto x - f(x) \cdot 3e_1$ envoie $\Delta$ dans $\Delta$. Alors $f(e_1 + e_2) = 2/3$. Une telle $f$ existe : $f(x_1, x_2) = 2x_1/3$.
\end{enumerate}
Le nombre de Cartan $\langle\gamma_{3e_1}, e_1 + e_2\rangle$ vaut donc $2/3$, qui n'est pas entier, et $3e_1 = 3 \cdot e_1$ est divisible dans $P$.
\end{contrexemple}

\begin{proof}
On a $e_1 = 3e_1 - (e_1 + e_2) - (e_1 - e_2)$ et $e_2 = (e_1 + e_2) - e_1$, d'où l'axiome a). Pour b), deux racines proportionnelles sont sur une même droite par l'origine, et chacune des quatre droites $\mathbb Q e_1$, $\mathbb Q e_2$, $\mathbb Q(e_1 + e_2)$, $\mathbb Q(e_1 - e_2)$ ne contient que deux racines opposées ; on vérifie les soixante-quatre couples. Pour c), les quatre réflexions du groupe de Weyl de type $B_2$, $(x_1, x_2) \mapsto (-x_1, x_2)$, $(x_1, -x_2)$, $(-x_2, -x_1)$ et $(x_2, x_1)$, permutent $\Delta$ et changent en leur opposé respectivement $\pm 3e_1$, $\pm 3e_2$, $\pm(e_1 + e_2)$ et $\pm(e_1 - e_2)$.

Pour 4), l'additivité donne $3 f(2e_1) = 2 f(3e_1) = 4$, donc $f(e_1 + e_2) + f(e_1 - e_2) = f(2e_1) = 4/3$. L'image de $e_1 + e_2$ est $(1 - 3f(e_1 + e_2),\, 1)$ ; c'est une racine de seconde coordonnée~$1$, donc $e_1 + e_2$ ou $-e_1 + e_2$, et $f(e_1 + e_2) \in \{0, 2/3\}$. De même l'image $(1 - 3f(e_1 - e_2),\, -1)$ de $e_1 - e_2$ donne $f(e_1 - e_2) \in \{0, 2/3\}$. La somme valant $4/3$, les deux valent $2/3$. Pour $f(x_1, x_2) = 2x_1/3$, l'application $x \mapsto x - f(x) \cdot 3e_1$ est la première réflexion. Enfin $3e_1 = 3 \cdot e_1$ avec $e_1 \in P$ et $3 \neq \pm 1$.
\end{proof}

\begin{remarque}
La réserve de la piste se confirme : les axiomes a) à c), même avec $n \in \mathbb Q$ dans b), ne donnent pas l'intégralité, et la symétrie de $3e_1$ y est unique, de sorte que le nombre $2/3$ ne dépend d'aucun choix. Les longueurs des racines sont dans le rapport $3/\sqrt 2$ au lieu de $\sqrt 2$ : le groupe engendré par les symétries est celui de $B_2$, mais $\Delta$ n'est pas cristallographique. La fiche ajoutait que $3e_1$ appartient encore au système simple $\{3e_1, e_2 - e_1\}$ ; c'est vrai au sens des groupes de réflexions, à coefficients rationnels positifs, mais pas au sens entier, car $e_1 + e_2 = \frac23 \cdot 3e_1 + (e_2 - e_1)$ n'est pas combinaison entière de $3e_1$ et $e_2 - e_1$ (vérifié dans le fichier). Un élément d'un système fondamental au sens entier fait partie d'une base de $P$, donc il est indivisible. Le crochet de la page n'est donc pas réfuté : les trois axiomes ne fournissent pas l'existence d'un tel système. Reste à lire sur le fac-similé le mot qui gouverne l'assertion, pour savoir si la page affirme, se propose de démontrer ou demande.
\end{remarque}

Ne sont pas démontrés ici : l'unicité des symétries par la forme invariante $\sum_{a} \langle x, a\rangle\langle y, a\rangle$ dans le cas général ; l'énoncé qu'un élément d'un système fondamental au sens entier est indivisible ; les axiomes de la donnée radicielle $(P, P', \Delta, \Delta', \varphi)$ de la seconde moitié du feuillet.
